\section{\href{https://doi.org/10.1088/1361-648X/adbfec}{Symmetry-adapted models for the hyperpolarizability of water}}

\subsection*{Supervisor: Dr David Wilkins}



Accurately modeling nonlinear optical experiments such as second-harmonic 
scattering and hyper-Raman spectroscopy requires the hyperpolarizability 
$\beta$, a nonlinear response to an applied electric field. The 
hyperpolarizability tensor is a computationally expensive quantity to 
calculate, making it a natural target for machine-learning methods. We test 
a family of recently developed models for the hyperpolarizability of water, 
trained on small clusters containing up to 8 water molecules. These models 
are able to predict $\beta$ for larger clusters, with more complex structures 
than those observed in the training set. For configurations of bulk water, 
the agreement is not so straightforward: while the total hyperpolarizability 
is quite well described, the predicted molecular $\beta$ tensors vary wildly 
between models. This means that while experiments whose outputs depend on 
total hyperpolarizability can be accurately modeled, those that require 
molecular quantities will require improved models.
